\documentclass[12pt]{jlreq}
\usepackage{amsmath, amssymb}

\newcommand{\C}{\mathbb{C}}
\newcommand{\R}{\mathbb{R}}
\newcommand{\Z}{\mathbb{Z}}
\newcommand{\N}{\mathbb{N}}
\newcommand{\F}{\mathbb{F}}
\newcommand{\Prob}{\mathbb{P}}
\newcommand{\Exp}{\mathbb{E}}
\newcommand{\dd}{\,d}
\newcommand{\Ker}{\operatorname{Ker}}
\newcommand{\Impart}{\operatorname{Im}}
\newcommand{\Hom}{\operatorname{Hom}}

\begin{document}

\(D\) を \(\R^n\) の開集合（ただし \(n \ge 2\)）とし，実数 \(1 \le p < \infty\) に対して
\[
  A^p := \{ u \in L^p(D) \mid u \text{ は } D \text{ 上}，\Delta u = 0 \text{ を超関数の意味で満たす。} \}
\]
と定める．ただし，\(\Delta u = \frac{\partial^2 u}{\partial x_1^2} + \dots + \frac{\partial^2 u}{\partial x_n^2}\) とし，偏微分は超関数の意味で考える．

(1) \(A^p\) 内の関数列 \(u_1, u_2, u_3, \dots\) が \(L^p(D)\) の位相で関数 \(u\) に強収束すれば，極限関数 \(u\) は \(A^p\) に属することを示せ．また，\(A^p\) 内の関数列 \(u_1, u_2, u_3, \dots\) が \(L^p(D)\) において関数 \(u\) に弱収束すれば，極限関数 \(u\) は \(A^p\) に属するか判定せよ．

(2) \(D\) が \(\R^2\) 上の単位円板の場合，\(A^1\) に属するが \(A^2\) に属さない関数の例を挙げよ．

\end{document}